\section{The Coupled Dynamics of p53 and DNA Across Binding Sequences}

\begin{figure}
\centering
\includegraphics[width=\textwidth]{Promoters/p53_Bend_vs_E2E.png}
\caption[BAX DNA Bend vs End-to-End Distance]{The DNA bend vs DNA end-to-end distance for the BAX promoter. p53-bound DNA is shown in blue, p53-free DNA is shown in red.}
\label{fig:promoters:bend-vs-e2e}
\end{figure}

\begin{figure}
\centering
\includegraphics[width=\textwidth]{Promoters/p53_Roll_vs_E2E.png}
\caption[BAX DNA Roll vs End-to-End Distance]{The average absolute values of DNA roll across nucleotides vs DNA end-to-end distance for the BAX promoter. p53-bound DNA is shown in blue, p53-free DNA is shown in red.}
\label{fig:promoters:roll-vs-e2e}
\end{figure}

\begin{figure}
\centering
\includegraphics[width=\textwidth]{Promoters/DNA_Buckle_compare.png}
\caption[p53 Response Element \(\Delta\)Buckle]{The \(\Delta\)Buckle, computed as (p53-free DNA Buckle) \textminus{} (p53-bound DNA Buckle). A value greater than zero means the p53-free DNA has greater buckle (bending of bases like a hinge) than the p53-bound. p53-interacting nucleotides are indicated with a box.}
\label{fig:promoters-dna-buckle-compare}
\end{figure}

\begin{figure}
\centering
\includegraphics[width=\textwidth]{Promoters/DNA_Propeller_compare.png}
\caption[p53 Response Element \(\Delta\)Propeller]{The \(\Delta\)Propeller, computed as (p53-free DNA Propeller) \textminus{} (p53-bound DNA Propeller). A value greater than zero means the p53-free DNA has greater propeller (twisting of two bases relative to one another) than the p53-bound. p53-interacting nucleotides are indicated with a box.}
\label{fig:promoters-dna-propeller-compare}
\end{figure}

\begin{figure}
\centering
\includegraphics[width=\textwidth]{Promoters/DNA_Tip_compare.png}
\caption[p53 Response Element \(\Delta\)Tip]{The \(\Delta\)Tip, computed as (p53-free DNA Tip) \textminus{} (p53-bound DNA Tip). A value greater than zero means the p53-free DNA has greater tip (rotation of a base pair around its short axis) than the p53-bound. p53-interacting nucleotides are indicated with a box.}
\label{fig:promoters-dna-tip-compare}
\end{figure}

\begin{figure}
\centering
\includegraphics[width=\textwidth]{Promoters/DNA_incline_compare.png}
\caption[p53 Response Element \(\Delta\)Inclination]{The \(\Delta\)Inclination, computed as (p53-free DNA Inclination) \textminus{} (p53-bound DNA Inclination). A value greater than zero means the p53-free DNA has greater inclination (rotation of a base pair around its long axis) than the p53-bound. p53-interacting nucleotides are indicated with a box.}
\label{fig:promoters-dna-inclination-compare}
\end{figure}

\begin{figure}
\centering
\includegraphics[width=\textwidth]{Promoters/DNA_Roll_compare.png}
\caption[p53 Response Element \(\Delta\)Roll]{The \(\Delta\)Roll, computed as (p53-free DNA Roll) \textminus{} (p53-bound DNA Roll). A value greater than zero means the p53-free DNA has greater roll (rotation of two bases relative to one another along their short axis) than the p53-bound. p53-interacting nucleotides are indicated with a box.}
\label{fig:promoters-dna-roll-compare}
\end{figure}